\chapter[1975 --- Baltimore et al.]{1975: David Baltimore, Renato Dulbecco, Howard Martin Temin}
\label{ch:1975_Baltimore}

\section*{Main Contribution}

\textit{Discovered how tumor viruses interact with the genetic material of the cell, explaining how viruses can cause cancer by inserting their genes into host DNA.}

\section{Official Citation}

for their discoveries concerning the interaction between tumour viruses and the genetic material of the cell

The 1975 Nobel Prize in Physiology or Medicine was awarded jointly to David Baltimore, Renato Dulbecco, and Howard Martin Temin for their pioneering discoveries concerning the interaction between tumour viruses and the genetic material of the cell. Their collective work fundamentally transformed our understanding of how viruses cause cancer and, in doing so, laid the foundations for modern molecular oncology, retrovirology, and gene therapy. Each laureate approached the problem from a different angle---Baltimore through biochemistry, Temin through virology and cell biology, and Dulbecco through tumor biology---yet their converging lines of evidence produced one of the most significant insights in the history of biology: that genetic information can flow not only from DNA to RNA to protein, but also from RNA back to DNA.

\section{Background and Historical Context}

The mid-twentieth century was a period of notable ferment in the biological sciences. The double-helical structure of DNA had been elucidated by Watson and Crick in 1953, and the central dogma of molecular biology---that genetic information flows from DNA to RNA to protein---was being elaborated by Francis Crick and others throughout the late 1950s and 1960s. At the same time, a parallel revolution was underway in cancer research. The discovery of tumor viruses---agents that could transform normal cells into malignant ones---had opened an entirely new avenue for understanding the origins of cancer.

The study of tumor viruses had a distinguished pedigree. Peyton Rous had identified a virus capable of inducing sarcomas in chickens as early as 1911, though his discovery was largely ignored for decades. In the 1950s and 1960s, Ludwig Gross, Sarah Stewart, and others isolated polyomavirus and other agents capable of causing tumors in rodents. By the mid-1960s, a critical distinction had emerged between two classes of tumor viruses: DNA tumor viruses, such as polyomavirus and simian virus 40 (SV40), which contain DNA as their genetic material, and RNA tumor viruses (later called retroviruses), which contain RNA. The RNA tumor viruses were of particular interest because they were known to cause leukemias and sarcomas in animals, and there was growing suspicion that analogous agents might play a role in human cancers.

Renato Dulbecco, an Italian-born virologist working at the Salk Institute in La Jolla, California, had been studying the mechanisms by which DNA tumor viruses transform cells. His work with polyomavirus and SV40 had established that these viruses integrate their DNA into the host cell genome, and that this integration event was closely associated with the malignant transformation of the cell. Dulbecco's laboratory became one of the preeminent centers for the study of tumor viruses, and his approach---combining rigorous quantitative biology with elegant experimental design---influenced an entire generation of virologists, including both Baltimore and Temin, who trained in or visited his laboratory.

David Baltimore, born in New York City in 1938, was a skilled young biochemist who had trained with James Darnell and Salvatore Luria at the Massachusetts Institute of Technology (MIT). Baltimore's early work focused on the replication of RNA viruses, and he had developed sophisticated biochemical tools for studying RNA synthesis in virus-infected cells. His work on vesicular stomatitis virus (VSV), an RNA virus, had already established him as one of the leading figures in molecular virology.

Howard Martin Temin, born in Philadelphia in 1934, was a virologist at the University of Wisconsin who had been studying Rous sarcoma virus (RSV) since his graduate days with Harry Rubin in the early 1960s. Temin had developed a cell culture system in which RSV could be studied quantitatively, and his observations of the virus's replication cycle led him to propose a radical hypothesis that would define his career.

The scientific milieu of the late 1960s and early 1970s was charged with controversy. The prevailing orthodoxy of the central dogma held that the flow of genetic information was strictly unidirectional: DNA $\rightarrow$ RNA $\rightarrow$ protein. The idea that information could flow in the reverse direction---from RNA to DNA---was considered heretical by many molecular biologists, including Francis Crick himself. This was the intellectual backdrop against which Temin and Baltimore made their notable discoveries.

\section{The Discovery/Contribution}

\subsection{Howard Temin and the Provirhus Hypothesis}

Temin's path to the discovery of reverse transcriptase began with a series of meticulous observations about the replication cycle of Rous sarcoma virus. Working in the late 1950s and 1960s, Temin noticed that when RSV infected chicken embryo fibroblasts, the virus did not simply produce new virus particles immediately. Instead, there was a ``latent period'' during which the infected cells became permanently altered---they grew faster, lost contact inhibition, and acquired other properties of transformed cells. Moreover, this transformed state was heritable: the daughter cells of virus-infected cells remained transformed even after the infectious virus had been removed.

These observations led Temin to propose, in 1964, the ``provirus hypothesis.'' Temin suggested that RSV copied its RNA genome into a DNA form that became integrated into the host cell's chromosomal DNA. This DNA provirus, Temin argued, served as the template for the production of new viral RNA and new virus particles. In essence, Temin proposed that information could flow from RNA to DNA---a direct contradiction of the central dogma as it was then understood.

The provirus hypothesis was met with widespread skepticism and, in some cases, outright hostility. Many leading molecular biologists regarded it as impossible or, at best, poorly supported by evidence. Temin himself later recalled that his hypothesis was ``a heresy'' that made him something of a pariah at scientific meetings. Nevertheless, Temin persevered, designing increasingly clever experiments to test his idea. He showed that treatment of virus-infected cells with actinomycin D, an inhibitor of DNA-dependent RNA synthesis, blocked the production of new virus---a result consistent with the idea that a DNA intermediate was required. He also demonstrated that infection of cells with RSV led to the formation of a DNA copy of the viral genome, using hybridization experiments to detect the proviral DNA.

\subsection{David Baltimore and the Discovery of Reverse Transcriptase}

In the spring of 1970, David Baltimore and his postdoctoral fellow H. M. Temin (working independently but arriving at the same conclusion almost simultaneously) announced one of the major discoveries in the history of molecular biology: the enzyme that copies RNA into DNA. Working with Rauscher murine leukemia virus (an RNA tumor virus closely related to RSV), Baltimore developed a biochemical assay for RNA-directed DNA polymerase activity. He mixed disrupted virus particles with radioactively labeled nucleotides and found that the virus itself contained an enzyme capable of synthesizing DNA using an RNA template.

The key to Baltimore's discovery was his insight that if the provirus hypothesis was correct, then the virus particle itself should carry an enzyme---an RNA-directed DNA polymerase---that could initiate the process of converting the viral RNA genome into DNA. To test this, Baltimore designed an elegant in vitro assay. He disrupted the viral particles and provided them with synthetic RNA templates, DNA primers, and radioactively labeled deoxyribonucleoside triphosphates. Within minutes, he observed robust DNA synthesis that was dependent on the presence of RNA template. The enzyme was later named reverse transcriptase, reflecting its ability to catalyze the reverse flow of genetic information from RNA to DNA.

Temin independently discovered the same enzyme in Rous sarcoma virus using a slightly different experimental approach. Temin and his student Satoshi Mizutama published their results in the same issue of \emph{Nature} as Baltimore, in June 1970. The simultaneous discovery by two independent groups provided powerful confirmation of the provirus hypothesis and established reverse transcriptase as a central enzyme in the biology of retroviruses.

\subsection{Renato Dulbecco and Tumor Virus DNA Integration}

Dulbecco's contribution to this triumvirate of discoveries was somewhat different in character but no less important. While Baltimore and Temin were elucidating the enzymatic mechanism of reverse transcription, Dulbecco was studying the consequences of tumor virus DNA integration into the host cell genome. Using polyomavirus and SV40 as model systems, Dulbecco demonstrated that the amount of viral DNA that became integrated into the host genome was directly correlated with the degree of cellular transformation. He showed that only a small fragment of the viral genome---the ``transforming region''---was sufficient to convert a normal cell into a cancer cell, while the remainder of the viral genome carried the instructions for virus replication.

Dulbecco's work established several principles that were essential for understanding how tumor viruses cause cancer. First, he demonstrated that viral transformation was a consequence of the stable integration of viral genetic material into the host cell genome. Second, he showed that this integration event could alter the expression of host cell genes, including genes involved in cell growth and division. Third, his mapping of the transforming region of SV40 provided the conceptual framework for the later identification of viral oncogenes---genes encoded by tumor viruses that are responsible for their cancer-causing ability.

Dulbecco's quantitative approach to virology was also instrumental in establishing the field of tumor biology on a rigorous footing. His methods for measuring virus yields, transformation frequencies, and DNA synthesis rates became standard tools in laboratories around the world. Many of the scientists who would go on to make major contributions to molecular biology and cancer research---including Baltimore, Temin, Michael Bishop, Harold Varmus, and others---trained in Dulbecco's laboratory or were significantly influenced by his approach.

\subsection{The Unified Impact}

Together, the work of Baltimore, Dulbecco, and Temin established the field of retrovirology and provided the conceptual and experimental foundations for understanding how viruses cause cancer at the molecular level. The discovery of reverse transcriptase was a paradigm-shifting event. It not only validated the provirus hypothesis and explained how RNA tumor viruses replicate, but it also demonstrated that the central dogma of molecular biology was not as inviolable as had been assumed. Information could flow from RNA to DNA, opening up entirely new possibilities for understanding gene expression and genome organization.

Moreover, the study of tumor virus DNA integration led directly to the discovery of oncogenes and proto-oncogenes---the normal cellular genes whose dysregulation can cause cancer. This line of investigation, pursued by Bishop, Varmus, and many others in the years following the 1975 Nobel Prize, would transform our understanding of the molecular basis of cancer and lead to the development of targeted cancer therapies.

\section{Impact on Modern Medicine}

The discoveries celebrated by the 1975 Nobel Prize have had an enormous and lasting impact on medicine and biology. Perhaps the most immediate practical consequence was the development of reverse transcriptase as a tool for molecular biology. The enzyme became indispensable for the construction of complementary DNA (cDNA) libraries, for mapping gene expression, and for cloning genes from their messenger RNA. Virtually every advance in molecular biology over the past five decades has depended, at some level, on the availability of reverse transcriptase.

In the field of virology, the understanding of retroviral replication that emerged from the work of Baltimore, Dulbecco, and Temin proved to be of paramount importance when the human immunodeficiency virus (HIV) was identified as the cause of AIDS in the early 1980s. HIV is a retrovirus, and its life cycle depends on reverse transcriptase. The development of antiretroviral drugs---the ``cocktail'' therapies that have transformed HIV from a death sentence into a manageable chronic condition---was based directly on the understanding of reverse transcriptase that originated with the 1975 Nobel discoveries. Drugs such as zidovudine (AZT), lamivudine, and tenofovir all target reverse transcriptase, inhibiting its activity and thereby preventing the virus from replicating. The success of these drugs stands as one of the most dramatic demonstrations of how basic scientific research can translate into life-saving medical treatments.

The concept that viral DNA integration into the host genome can cause cancer---a principle established by Dulbecco and his colleagues---has also proved to be of enduring importance. The discovery that retroviruses can activate or disrupt cellular oncogenes through insertional mutagenesis led directly to the identification of numerous human oncogenes and tumor suppressor genes. This knowledge has, in turn, enabled the development of targeted cancer therapies, such as imatinib (Gleevec) for chronic myeloid leukemia and trastuzumab (Herceptin) for HER2-positive breast cancer. The principle that cancer is fundamentally a disease of the genome---a principle that emerged from the study of tumor viruses---now underpins the entire field of precision oncology.

Reverse transcriptase has also become a cornerstone tool in genomics and transcriptomics. The technique of quantitative reverse transcription polymerase chain reaction (RT-qPCR) is the gold standard for measuring gene expression and is used daily in research laboratories and clinical diagnostic settings worldwide. The development of RNA sequencing technologies, which underlie much of modern genomics, also depends on reverse transcription of RNA into DNA. Without the enzyme discovered by Baltimore and Temin, the Human Genome Project, the ENCODE project, and the revolution in single-cell transcriptomics would not have been possible.

In the realm of gene therapy, the understanding of how retroviruses integrate their genetic material into host cell chromosomes has been exploited to develop viral vectors for gene delivery. Retroviral vectors---based on the same principles elucidated by Dulbecco, Baltimore, and Temin---have been used to correct genetic defects in patients with severe combined immunodeficiency (SCID), sickle cell disease, and other inherited disorders. Although early gene therapy trials using retroviral vectors sometimes caused insertional mutagenesis and leukemia, more recent approaches using lentiviral vectors (derived from HIV) have proved both safer and more effective, and several gene therapies based on these principles have now been approved for clinical use.

The work of the 1975 laureates also had a significant impact on our understanding of evolution and genome biology. The recognition that reverse transcriptase can generate DNA copies of RNA molecules has revealed the existence of retroelements---mobile genetic elements that transpose via an RNA intermediate---that make up a substantial fraction of the human genome. Approximately 40\% of the human genome consists of sequences derived from retrotransposons, including endogenous retroviruses, long interspersed nuclear elements (LINEs), and short interspersed nuclear elements (SINEs). The study of these elements has become a major field in its own right, with implications for understanding genome evolution, gene regulation, and genetic disease.

\section{Legacy}

The legacy of Baltimore, Dulbecco, and Temin is woven into the fabric of modern biology and medicine. David Baltimore went on to become one of an influential figures in American science, serving as president of the California Institute of Technology and founding the Whitehead Institute for Biomedical Research at MIT. His later work on the molecular biology of HIV, the regulation of gene expression, and the interferon response system further cemented his reputation as one of the great biologists of the twentieth century. Baltimore's insistence on rigorous experimental standards and his ability to bring biochemical approaches to bear on complex biological problems influenced generations of scientists.

Howard Temin remained at the University of Wisconsin-Madison for his entire career, continuing his studies of retroviruses and their role in cancer. He became a passionate advocate for cancer prevention and public health, arguing that understanding the viral causes of cancer should translate into efforts to prevent viral infections through vaccination and other public health measures. Temin was also outspoken in his defense of basic research, arguing that a major medical advances often come from investigations motivated purely by curiosity rather than by immediate practical goals.

Renato Dulbecco continued his work at the Salk Institute, turning his attention to the molecular biology of cancer and, later in his career, to the Human Genome Project. Dulbecco's 1986 editorial in \emph{Science}, in which he argued that understanding cancer required a complete knowledge of the human genome, is widely credited with helping to galvanize support for the Human Genome Project, which was launched in 1990 and completed in 2003. The Human Genome Project, in turn, has transformed every aspect of biology and medicine, from the diagnosis and treatment of genetic diseases to the understanding of human evolution and diversity.

The three laureates of 1975 demonstrated that fundamental scientific discoveries---made without regard to immediate practical applications---can have consequences that are both significant and far-reaching. Their work on tumor viruses and the interaction between viral and cellular genetic material opened up entirely new fields of inquiry and provided tools and concepts that continue to shape medicine and biology to this day. The discovery of reverse transcriptase, in particular, stands as one of the major achievements in the history of the life sciences---a discovery that overturned established dogma, created new possibilities for research and therapy, and earned its discoverers the highest honor in medicine.


\section{Modern Developments}

Baltimore, Dulbecco, and Temin's work on tumor viruses has led to modern cancer genomics. Understanding of reverse transcriptase has enabled HIV treatment (NRTIs) and mRNA vaccine development. Understanding of viral oncogenes has revealed cancer signaling pathways, leading to targeted therapies (imatinib, erlotinib). Understanding of retroviral integration has enabled gene therapy vectors. Understanding of host-virus interactions has informed development of antiviral drugs.
